\BeleskeID{09_2-s018}
% element: 09_2-s018:e04
Važno je da se u izrazu za prag javlja odnos parametara P i N tranzistora. Isti izraz se može preurediti tako da pod korenom bude normalizovan odnos:
\[\boldsymbol{\begin{aligned}V_{IL}&=V_{Tn}+(V_{DD}+V_{Tp})\frac{k_p}{k_n}\sqrt{\frac{k_p}{k_p+k_n}}\\&=V_{Tn}+(V_{DD}+V_{Tp})\frac{k_p}{k_n}\sqrt{\frac{\frac{k_p}{k_n}}{1+\frac{k_p}{k_n}}}\end{aligned}}\]
Ako se obe širine povećaju istim faktorom $F$, taj faktor se skrati u odnosu $k_p/k_n$. Zato se u ovoj statičkoj analizi pragovi ne menjaju. Isto zaključivanje važi i na drugim mestima gde izrazi zavise samo od odnosa ovih parametara, uz nepromenjene ostale pretpostavke.
